% !TEX program = xelatex
\documentclass[14pt]{extarticle}

\usepackage[
  a4paper,
  left=2.5cm,
  right=1.5cm,
  top=2cm,
  bottom=2cm
]{geometry}

\usepackage[ukrainian]{babel}
\usepackage{fontspec}
\setmainfont{Times New Roman}

\usepackage{setspace}
\onehalfspacing
\setlength{\parindent}{1.27cm}
\setlength{\parskip}{0pt}
\usepackage{microtype}

\usepackage{titlesec}
\titleformat{\section}
  {\bfseries\fontsize{16pt}{24pt}\selectfont\centering}
  {\thesection}{0.5em}{}
\titlespacing*{\section}{0pt}{20pt}{6pt}

\titleformat{\subsection}
  {\bfseries\fontsize{14pt}{21pt}\selectfont}
  {\thesubsection}{0.5em}{}
\titlespacing*{\subsection}{1.27cm}{6pt}{6pt}

\usepackage{enumitem}
\setlist[itemize,1]{
  label=$\bullet$,
  leftmargin=1.27cm,
  labelwidth=0.5cm,
  labelsep=0.3cm,
  itemindent=0pt,
  itemsep=3pt,
  topsep=6pt,
  parsep=0pt,
  partopsep=0pt
}
\setlist[enumerate,1]{
  label=\arabic*.,
  leftmargin=2.54cm,
  labelwidth=0.63cm,
  labelsep=0.54cm,
  itemindent=0pt,
  itemsep=0pt,
  topsep=0pt,
  parsep=0pt,
  partopsep=0pt
}

\usepackage{array}
\usepackage{longtable}
\usepackage{booktabs}
\usepackage{tabularx}
\usepackage{multirow}
\usepackage{graphicx}

\usepackage[hidelinks,breaklinks=true]{hyperref}

\begin{document}
\sloppy
\emergencystretch=3em

% ==================== ТИТУЛ ====================
\begin{center}
\textbf{Лабораторна робота №3}

\medskip

\textbf{\MakeUppercase{Визначення стейкхолдерів та складності реалізації проєкту з точки зору його управління. Формування вимог до проєкту.}}
\end{center}

\medskip

\textbf{Мета роботи:} навчитись визначати всі зацікавлені сторони проєкту та їх вимоги. Навчитись планувати виконання проєкту відповідно до вимог зацікавлених сторін та формувати матрицю відповідальності (RAM). Визначити складність проєкту з точки зору управління.

\textbf{Проєкт:} <<Система для автоматичної нормалізації суржику та діалектних форм української мови>>.

% ==================== РОЗДІЛ 1: СТЕЙКХОЛДЕРИ ====================
\section*{1. Визначення стейкхолдерів та їх вплив на проєкт}
\addcontentsline{toc}{section}{1. Визначення стейкхолдерів та їх вплив на проєкт}

Стейкхолдером є будь-хто, хто впливає на проєкт або на кого впливає проєкт -- позитивно чи негативно. Нижче наведено повний перелік зацікавлених сторін з оцінкою їх потужності (Power) та зацікавленості (Interest).

\begin{longtable}{|p{0.04\textwidth}|p{0.29\textwidth}|p{0.17\textwidth}|p{0.09\textwidth}|p{0.10\textwidth}|p{0.14\textwidth}|}
\hline
\textbf{ID} & \textbf{Стейкхолдер} & \textbf{Роль у проєкті} & \textbf{Power} & \textbf{Interest} & \textbf{За / Проти} \\
\hline
\endfirsthead
\hline
\textbf{ID} & \textbf{Стейкхолдер} & \textbf{Роль у проєкті} & \textbf{Power} & \textbf{Interest} & \textbf{За / Проти} \\
\hline
\endhead
S1 & Коваль Денис Петрович & Виконавець, менеджер проєкту & Висока & Висока & За \\
\hline
S2 & Сивоконь Олексій Олегович & Науковий керівник & Висока & Висока & За \\
\hline
S3 & Кафедра СШІ, НУ <<Львівська Політехніка>> & Інституційний замовник & Висока & Середня & За \\
\hline
S4 & Команда Lapa LLM (Laba.ai) & Розробники базової моделі & Низька & Середня & За \\
\hline
S5 & Носії діалектів (гуцульського, бойківського, закарпатського) & Кінцеві користувачі -- джерело задачі & Низька & Висока & За \\
\hline
S6 & Внутрішні мігранти та переселенці (суржикомовці) & Кінцеві користувачі & Низька & Висока & За \\
\hline
S7 & Ukrainian NLP-спільнота (дослідники, автори Vuyko Mistral, UA-GEC) & Академічна аудиторія, рецензенти & Середня & Висока & За \\
\hline
S8 & EdTech-платформи (Prometheus, МОН) & Потенційні інтегратори результату & Середня & Середня & Нейтрально \\
\hline
S9 & Мовні активісти / ТМШ & Промоутери норми, потенційні критики & Низька & Висока & За \\
\hline
S10 & NLP-розробники у tech-компаніях & Потенційні споживачі API/компонента & Середня & Середня & За \\
\hline
\end{longtable}

% ==================== РОЗДІЛ 2: СІТКА ОЦІНКИ ====================
\section*{2. Сітка оцінки зацікавлених сторін (Power/Interest)}
\addcontentsline{toc}{section}{2. Сітка оцінки зацікавлених сторін}

Сітка оцінки розміщує стейкхолдерів у чотирьох квадрантах за осями «Потужність» (Power) та «Зацікавленість» (Interest). Це визначає стратегію взаємодії з кожним стейкхолдером.

\begin{center}
\begin{tabular}{|p{0.01\textwidth}|p{0.42\textwidth}|p{0.42\textwidth}|}
\hline
& \multicolumn{1}{c|}{\textbf{Interest: Низька}} & \multicolumn{1}{c|}{\textbf{Interest: Висока}} \\
\hline
\rotatebox[origin=c]{90}{\textbf{Power: Висока}} &
\textbf{Квадрант II -- Задовольняти} \newline S3 Кафедра СШІ / НУ «ЛП» &
\textbf{Квадрант I -- Ключові гравці} \newline S1 Коваль Денис \newline S2 Сивоконь О.О. \\
\hline
\rotatebox[origin=c]{90}{\textbf{Power: Низька}} &
\textbf{Квадрант III -- Моніторинг} \newline S4 Команда Lapa LLM &
\textbf{Квадрант IV -- Інформувати} \newline S5 Носії діалектів \newline S6 Мігранти/переселенці \newline S7 Ukrainian NLP-спільнота \newline S9 Мовні активісти / ТМШ \\
\hline
\end{tabular}
\end{center}

\medskip

Стейкхолдери S8 (EdTech) та S10 (NLP-розробники) розміщуються між квадрантами II і IV (середній рівень за обома осями) -- їх варто залучати ситуативно.

\textbf{Стратегії взаємодії:}
\begin{itemize}
  \item \textbf{Квадрант I (Ключові гравці, High Power / High Interest):} активне залучення на всіх фазах, регулярна комунікація. S1 і S2 -- основна команда прийняття рішень.
  \item \textbf{Квадрант II (Задовольняти, High Power / Low Interest):} своєчасне інформування про прогрес, дотримання формальних вимог кафедри. S3 погоджує тему, приймає результат.
  \item \textbf{Квадрант III (Моніторинг, Low Power / Low Interest):} пасивний моніторинг. S4 -- їх відкрита модель використовується, достатньо публічного звіту.
  \item \textbf{Квадрант IV (Інформувати, Low Power / High Interest):} регулярне публічне інформування через публікацію датасетів, моделей, демо. S5, S6, S9 -- представники цільової аудиторії; S7 -- академічна спільнота отримує доступ до репозиторію та результатів.
\end{itemize}

% ==================== РОЗДІЛ 3: ВИМОГИ СТЕЙКХОЛДЕРІВ ====================
\section*{3. Вимоги стейкхолдерів та фази їх виконання}
\addcontentsline{toc}{section}{3. Вимоги стейкхолдерів та фази їх виконання}

\begin{longtable}{|p{0.03\textwidth}|p{0.25\textwidth}|p{0.38\textwidth}|p{0.18\textwidth}|}
\hline
\textbf{ID} & \textbf{Стейкхолдер} & \textbf{Вимоги та очікування} & \textbf{Фаза} \\
\hline
\endfirsthead
\hline
\textbf{ID} & \textbf{Стейкхолдер} & \textbf{Вимоги та очікування} & \textbf{Фаза} \\
\hline
\endhead
S1 & Виконавець & Чітка постановка задачі, доступ до ресурсів, відтворювані результати & Всі \\
\hline
S2 & Науковий керівник & Дотримання академічних стандартів, новизна, відповідність темі БКР & Ініціювання, Контроль, Закриття \\
\hline
S3 & Кафедра СШІ & Вчасна здача роботи, відповідність формальним вимогам та структурі БКР & Ініціювання, Закриття \\
\hline
S4 & Команда Lapa LLM & Коректне цитування моделі, можливий зворотній зв'язок про результати донавчання & Виконання, Закриття \\
\hline
S5 & Носії діалектів & Система коректно нормалізує їхній діалект; не спотворює зміст & Виконання, Контроль \\
\hline
S6 & Мігранти / переселенці & Простий та доступний веб-інтерфейс; підтримка суржику & Виконання, Контроль \\
\hline
S7 & Ukrainian NLP-спільнота & Відкритий код і датасети, описана методологія, відтворювані метрики & Виконання, Закриття \\
\hline
S8 & EdTech-платформи & Можливість інтеграції через API або як модуль; документація & Закриття \\
\hline
S9 & Мовні активісти & Система нормалізує до літературної норми, а не «русифікує» & Виконання, Контроль \\
\hline
S10 & NLP-розробники & Документований API або Gradio-інтерфейс, відкрита модель на HuggingFace & Закриття \\
\hline
\end{longtable}

% ==================== РОЗДІЛ 4: ФАЗИ ПРОЄКТУ ====================
\section*{4. Фази проєкту з урахуванням стейкхолдерів}
\addcontentsline{toc}{section}{4. Фази проєкту з урахуванням стейкхолдерів}

\textbf{Фаза 1. Ініціювання.} Основні учасники: S1, S2, S3. Результати: затверджена тема БКР, визначені стейкхолдери, узгоджена SMART-мета, документ сфери застосування. Замовник (S3 через S2) підписує погодження теми.

\textbf{Фаза 2. Планування.} Основні учасники: S1, S2; залучаються S7 (огляд літератури з їхніх робіт). Результати: план-графік, реєстр ризиків, вибір архітектури, визначені метрики. Враховуються вимоги S7 (відтворюваність), S5/S6 (якість нормалізації), S9 (мовна коректність).

\textbf{Фаза 3. Виконання.} Основні учасники: S1 (розробка). Залучаються: S4 (Lapa LLM як ресурс), S5/S6 (тестові приклади, оцінка якості), S7 (публікація проміжних результатів). Результати: навчена baseline-модель, мультидіалектний корпус, донавчана LLM, Gradio-інтерфейс.

\textbf{Фаза 4. Контроль.} Основні учасники: S1, S2. Залучаються S5/S6/S9 (оцінка якості нормалізації на реальних прикладах). Відстеження метрик (BLEU, chrF++, TER, COMET), консультації з керівником, корекція підходу.

\textbf{Фаза 5. Закриття.} Основні учасники: S1, S2, S3. Залучаються S4, S7, S8, S10 (публікація, документація). Результати: текст БКР, публічний репозиторій, модель на HuggingFace, захист.

% ==================== РОЗДІЛ 5: RAM ====================
\section*{5. Матриця розподілу відповідальності (RAM)}
\addcontentsline{toc}{section}{5. Матриця розподілу відповідальності (RAM)}

Умовні позначення: \textbf{A} -- відповідальний (Accountable), \textbf{P} -- учасник (Participant), \textbf{O} -- необхідна думка (Opinion required), порожньо -- не залучений.

\begin{center}
\begin{longtable}{|p{0.27\textwidth}|c|c|c|c|c|c|c|c|c|c|}
\hline
\textbf{Завдання / Етап} & \textbf{S1} & \textbf{S2} & \textbf{S3} & \textbf{S4} & \textbf{S5} & \textbf{S6} & \textbf{S7} & \textbf{S8} & \textbf{S9} & \textbf{S10} \\
\hline
\endfirsthead
\hline
\textbf{Завдання / Етап} & \textbf{S1} & \textbf{S2} & \textbf{S3} & \textbf{S4} & \textbf{S5} & \textbf{S6} & \textbf{S7} & \textbf{S8} & \textbf{S9} & \textbf{S10} \\
\hline
\endhead
Визначення теми та мети & A & P & O & & & & & & & \\
\hline
Огляд літератури та джерел & A & O & & & & & P & & & \\
\hline
Аналіз наборів даних & A & O & & O & & & O & & & \\
\hline
Підготовка корпусу (гуцульський) & A & O & & & O & & & & & \\
\hline
Синтетичне розширення даних & A & O & & O & O & O & & & O & \\
\hline
Навчання baseline seq2seq & A & O & & & & & & & & \\
\hline
Формування мультидіалектного корпусу & A & O & & & O & O & O & & O & \\
\hline
Донавчання Lapa LLM (QLoRA) & A & O & & P & & & & & & \\
\hline
Оцінювання якості (метрики) & A & P & & & O & O & O & & O & \\
\hline
Розробка Gradio-інтерфейсу & A & O & & & & P & & O & & O \\
\hline
Написання тексту БКР & A & P & O & & & & & & & \\
\hline
Публікація коду та моделі & A & O & & O & & & P & O & & O \\
\hline
Захист БКР & A & P & A & & & & & & & \\
\hline
\end{longtable}
\end{center}

% ==================== РОЗДІЛ 6: ВИМОГИ ТА СКЛАДНІСТЬ ====================
\section*{6. Вимоги до проєкту та складність з точки зору управління}
\addcontentsline{toc}{section}{6. Вимоги до проєкту та складність}

\subsection*{6.1. Функціональні вимоги}

Функціональні вимоги сформовані на основі потреб стейкхолдерів:

Система повинна приймати на вхід текст (до кількох речень) у форматі рядка та повертати нормалізований варіант літературною українською мовою. Підтримуються чотири режими вхідних даних: гуцульський діалект, бойківський діалект, закарпатський діалект та суржик (вимога S5, S6). Система надає веб-інтерфейс із вибором режиму та текстовим полем введення (вимога S6, S10). Код, навчені моделі та датасети публікуються відкрито (вимога S7, S10). Метрики якості (BLEU, chrF++, TER, COMET) документуються для кожного підтримуваного різновиду (вимога S2, S7).

\subsection*{6.2. Нефункціональні вимоги}

Відтворюваність: усі етапи підготовки даних і навчання мають бути задокументовані та відтворювані (вимога S7). Мовна коректність: нормалізація приводить до літературної норми без зміни змісту та без суржикових або русизмів у виведенні (вимога S9). Задокументованість: текст БКР відповідає вимогам кафедри та структурі наукової роботи (вимога S3). Доступність: веб-інтерфейс доступний без встановлення додаткового ПЗ (вимога S6).

\subsection*{6.3. Тип проєкту та складність}

Відповідно до матриці типів проєктів (стабільний/еволюційний $\times$ низький/високий виклик), проєкт відноситься до \textbf{Квадранту II -- Еволюційний / Середній виклик}:

\textbf{Еволюційні вимоги.} Задача нормалізації діалектів і суржику є новою та недостатньо формалізованою. На початку проєкту неможливо точно передбачити, яка якість буде досягнута на кожному діалекті та яким чином краще формувати синтетичні дані. Вимоги уточнюються по ходу: після baseline-етапу стало зрозумілим, що для мультидіалектного корпусу потрібне більш ретельне балансування.

\textbf{Рівень виклику -- Середній (між LOW та HIGH).} Більшість ключових компетенцій (Python, PyTorch, Hugging Face, NLP) присутні у виконавця. Проте донавчання великих мовних моделей (12B параметрів, QLoRA) та специфіка низькоресурсних діалектних корпусів виходять за межі стандартного університетського курсу -- потрібно залучати зовнішні ресурси (відкриті моделі, академічні статті, відкриті датасети).

\textbf{Критичність проєкту} оцінюється як \textbf{менеджментний рівень}: проєкт є ключовим для досягнення академічної мети (захист БКР), але не є критичним для виживання організації. Водночас результати мають практичну цінність для NLP-спільноти та потенційних інтеграторів (S7, S8, S10).

% ==================== ВИСНОВКИ ====================
\section*{Висновки}
\addcontentsline{toc}{section}{Висновки}

У ході виконання лабораторної роботи визначено 10 зацікавлених сторін проєкту -- від безпосередніх учасників (виконавець, науковий керівник, кафедра) до нетривіальних стейкхолдерів: команди Lapa LLM, носіїв діалектів, внутрішніх мігрантів, Ukrainian NLP-спільноти, EdTech-платформ, мовних активістів та NLP-розробників.

Побудовано сітку оцінки (Power/Interest) та визначено стратегію взаємодії з кожним квадрантом: ключові гравці (S1, S2) -- активне залучення на всіх фазах; кафедра (S3) -- своєчасне задоволення формальних вимог; NLP-спільнота та кінцеві користувачі -- публічне інформування через відкриті публікації.

Сформовано таблицю вимог стейкхолдерів, прив'язану до фаз проєкту, та матрицю розподілу відповідальності (RAM) для 13 ключових завдань. Визначено функціональні та нефункціональні вимоги до системи, а тип проєкту класифіковано як еволюційний з середнім рівнем виклику -- менеджментна критичність.

\end{document}
